\reportchapter{III. LÝ THUYẾT, GIẢI PHÁP, THUẬT TOÁN VÀ QUÁ TRÌNH PHÁT TRIỂN HỆ THỐNG}
\reportsection{a. Các lý thuyết, giải pháp và thuật toán liên quan}

\reportsubsection{Các lý thuyết nền tảng}
\textbf{DevOps:} là phương pháp kết nối hoạt động phát triển và vận hành thông qua một quy trình liên tục gồm viết mã, kiểm thử, đóng gói, triển khai và theo dõi hệ thống. Trong dự án, DevOps được thể hiện qua CI/CD, container, Infrastructure as Code, Helm và Kubernetes để giảm thao tác thủ công và bảo đảm mỗi phiên bản được triển khai theo cùng một quy trình \cite{devops-model}.

\textbf{DevSecOps} mở rộng DevOps bằng cách đưa bảo mật vào từng bước của quy trình phát triển, thay vì chỉ kiểm tra trước khi đưa hệ thống lên production. Các nội dung kiểm tra gồm mã nguồn, dependency, container image, secret, cấu hình hạ tầng và quyền truy cập. Với Philobiblus, các bước kiểm tra này được tích hợp vào CI/CD để phát hiện sớm rủi ro trước khi artifact được phát hành \cite{devsecops-model}.

\textbf{CI/CD} tự động hóa chuỗi kiểm thử, đóng gói và phát hành, giúp mỗi thay đổi được kiểm tra trước khi đi vào môi trường chạy \cite{devops-model,github-actions}.

\textbf{Infrastructure as Code} biểu diễn hạ tầng bằng code có thể được quản lý bằng Git, giúp quản lý hạ tầng nhất quán và cho phép xem trước thay đổi trước khi áp dụng \cite{iac-terraform,terraform-docs}.

\textbf{Containerization} giúp đóng gói ứng dụng cùng dependency trong image bất biến, giảm khác biệt giữa môi trường phát triển, kiểm thử và production \cite{docker}.

\textbf{Điều phối workload:} Kubernetes quản lý Pod, Deployment, Service, probe, rollout và autoscaling để workload có thể tự phục hồi và mở rộng khi lượng người dùng biến động \cite{kubernetes}.

\textbf{MLOps} là tập hợp các thực hành kết hợp giữa phát triển mô hình ML và vận hành hệ thống nhằm tự động hóa, chuẩn hóa và quản lý vòng đời mô hình một cách hiệu quả \cite{mlops-gcp}.

\textbf{Observability:} là việc đo lường và theo dõi trạng thái bên trong của một hệ thống phần mềm dựa trên các dữ liệu đầu ra mà nó tạo ra trong quá trình hoạt động \cite{observability-otel}.

\reportsubsection{Các giải pháp và công cụ áp dụng}
\textbf{Terraform:} là một công cụ giúp quản lý và tự động hóa hạ tầng công nghệ dưới dạng mã nguồn (IaC) thông qua các bước viết cấu hình, lập kế hoạch và áp dụng thay đổi \cite{terraform-docs}.

\textbf{Google Cloud và GKE:} GKE Autopilot cung cấp cluster Kubernetes managed, còn Google Cloud cung cấp các dịch vụ cloud để deploy hệ thống vào hoạt động \cite{gke-autopilot}.

\textbf{Docker} tạo image cho các thành phần của hệ thống như frontend, backend, mô hình gợi ý,... để có thể deploy lên các môi trường khác như cloud.

\textbf{Helm} giúp người dùng dễ dàng định nghĩa, cài đặt, nâng cấp và quản lý các cấu hình phức tạp trên cụm Kubernetes thông qua các gói đóng gói sẵn gọi là Chart \cite{helm-charts}.

\textbf{GitHub Actions:} là nền tảng tự động hóa và CI/CD tích hợp trực tiếp ngay trên GitHub \cite{github-actions}.

\textbf{DVC} (Data Version Control) là công cụ quản lý phiên bản dữ liệu và quy trình xử lý dữ liệu, hoạt động cùng với Git. \cite{dvc}.

\textbf{MLflow} là nền tảng theo dõi thí nghiệm và quản lý vòng đời mô hình học máy. Công cụ lưu thông tin về run, parameter, metric, artifact và model version \cite{mlflow}.

\textbf{Cloud SQL} là dịch vụ cơ sở dữ liệu quan hệ được quản lý hoàn toàn trên Google Cloud, trong dự án sử dụng PostgreSQL để lưu dữ liệu nghiệp vụ \cite{cloud-sql}.

\textbf{Redis} là máy chủ cấu trúc dữ liệu trong bộ nhớ, thường được dùng làm cache, queue hoặc lớp xử lý dữ liệu có yêu cầu phản hồi nhanh \cite{redis-datatypes}.

\textbf{Cloud Monitoring} là dịch vụ của Google Cloud dùng để thu thập, hiển thị và cảnh báo dựa trên các chỉ số vận hành của tài nguyên \cite{cloud-monitoring}.

\textbf{Managed Service for Prometheus:} Managed Service for Prometheus là dịch vụ quản lý Prometheus trên Google Cloud, dùng để thu thập và truy vấn metric do các workload Kubernetes xuất ra \cite{prometheus,managed-prometheus}.

\reportsubsection{Các thuật toán sử dụng}
\textbf{TF-IDF} biến metadata sách thành vector thưa, trong đó các từ xuất hiện nhiều trong một sách nhưng ít xuất hiện trong catalog có trọng số cao \cite{sklearn}.

\textbf{Cosine similarity} là phương pháp đo góc giữa hai vector TF-IDF để xếp hạng các sách có nội dung gần với sách nguồn mà không phụ thuộc độ dài tuyệt đối của chuỗi \cite{sklearn}.

\textbf{Snapshot hash và quality gate:} Hệ thống băm nội dung logic của snapshot, đối chiếu schema và mức suy giảm catalog trước khi cho phép candidate thay champion; cách kiểm soát này được triển khai riêng cho pipeline của Philobiblus và liên kết với dữ liệu, run cùng model version \cite{dvc,mlflow}.

\reportsection{b. Quá trình phát triển hệ thống}

\reportsubsection{Hoàn thiện mã nguồn và artifact}
\textbf{Frontend} được phát triển bằng React 18 và Vite. Giao diện tổ chức theo component-based architecture, sử dụng shadcn/ui với style base-nova, Tailwind CSS và Lucide icons; React Router quản lý điều hướng, còn các Context Provider tách trạng thái xác thực, chế độ xem và cài đặt. Các service module gọi REST API, nên frontend không thể truy cập trực tiếp vào database hoặc model artifact. Frontend được deploy trực tiếp lên GitHub Pages.

\textbf{Backend} dùng FastAPI và Uvicorn, tổ chức router theo domain như authentication, books, users, reviews và reading history. SQLAlchemy quản lý model và session database, Pydantic schema kiểm tra dữ liệu vào ra, còn dependency injection cung cấp database session cho từng request. JWT, quyền sở hữu và visibility đối với thông tin sách được kiểm tra khi gọi API; middleware bổ sung CORS, metric và trace. Cuối cùng, em đóng gói toàn bộ thành docker image để triển khai lên Kubernetes.

\textbf{Cơ sở dữ liệu và cache:} PostgreSQL lưu dữ liệu nghiệp vụ theo schema được khai báo bằng SQLAlchemy; môi trường production sử dụng Cloud SQL thay cho database container trong GKE. Redis đảm nhiệm cache catalog và rate limit, giúp giảm số lần truy vấn database và giới hạn các request tốn tài nguyên. Backend cấu hình connection pool, timeout và cơ chế bỏ qua cache khi Redis tạm thời không khả dụng.

\textbf{Recommendation service:} chạy bằng FastAPI, sử dụng các metadata của sách như tên, tác giả, thể loại để gợi ý sách tương đồng. Pipeline tạo snapshot catalog, chuẩn hóa dữ liệu và dùng TF-IDF để huấn luyện model Joblib; DVC quản lý các bước xử lý, còn MLflow lưu metric và phiên bản model. Khi hoạt động, hệ thống xếp hạng sách bằng cosine similarity và đưa ra top 5 sách có cosine similarity cao nhất.

\reportsubsection{Tạo Helm chart}
Hệ thống được tách thành 2 namespace: \texttt{philobiblus} chứa các workload phục vụ người dùng, còn \texttt{philobiblus-mlops} chứa các workload huấn luyện và phát hành model. Việc tách ra giúp hệ thống gợi ý sách có thể được phát triển và vận hành độc lập với các workload phục vụ người dùng, giảm rủi ro khi quá trình huấn luyện model gặp lỗi. Dưới đây là mô tả luồng hoạt động của hệ thống được thiết kế trong Helm chart:

\textbf{Luồng runtime:} Người dùng đi qua Gateway để vào nhóm workload ứng dụng. Gateway chuyển yêu cầu đến Backend; Backend đọc và ghi dữ liệu nghiệp vụ trong Cloud SQL, sử dụng Redis để lưu tạm dữ liệu thường dùng, rồi giao tiếp với Recommendation Service khi cần tạo gợi ý. Recommendation Service dùng model hiện hành đã được nạp trong Pod và trả kết quả về Backend để hoàn thành request.

\textbf{Luồng MLOps control:} Retrain CronJob trong namespace MLOps định kỳ lấy dữ liệu cần thiết, tạo snapshot, huấn luyện và đánh giá model. CronJob gửi run, metric và phiên bản model cho MLflow; trainer đồng thời kiểm tra Recommendation Service trước khi cho phép phát hành. Các workload gửi log và metric đến Cloud Monitoring để theo dõi rollout, lỗi và mức sử dụng tài nguyên. Các thao tác này được giới hạn bằng quyền của trainer, không làm namespace MLOps tham gia vào luồng request người dùng.

\textbf{Luồng artifact/model:} Sau khi model đạt điều kiện phát hành, trainer ghi model và thông tin release vào GCS. Recommendation Service lấy artifact tương ứng từ GCS, kiểm tra tính toàn vẹn rồi nạp model mới. Vì vậy, GCS đóng vai trò vùng lưu trữ trung gian giữa pipeline huấn luyện và workload phục vụ gợi ý.

\begin{figure}[H]
\centering
\definecolor{k8cleanInk}{HTML}{243447}
\definecolor{k8cleanCluster}{HTML}{64748B}
\definecolor{k8cleanAppStroke}{HTML}{2563EB}
\definecolor{k8cleanAppFill}{HTML}{EFF6FF}
\definecolor{k8cleanOpsStroke}{HTML}{7C3AED}
\definecolor{k8cleanOpsFill}{HTML}{F5F3FF}
\definecolor{k8cleanRuntime}{HTML}{1D4ED8}
\definecolor{k8cleanControl}{HTML}{C2410C}
\definecolor{k8cleanArtifact}{HTML}{15803D}
\definecolor{k8cleanExternalFill}{HTML}{F8FAFC}
\definecolor{k8cleanExternalStroke}{HTML}{475569}

\resizebox{\textwidth}{!}{%
\begin{tikzpicture}[
  x=1cm,y=1cm,
  >=latex,
  font=\scriptsize\sffamily,
  every node/.style={text=k8cleanInk},
  box/.style={draw=k8cleanInk, line width=0.75pt, rounded corners=3pt, align=center, inner sep=5pt, minimum height=9.2mm},
  app/.style={box, fill=k8cleanAppFill, draw=k8cleanAppStroke},
  ops/.style={box, fill=k8cleanOpsFill, draw=k8cleanOpsStroke},
  external/.style={box, fill=k8cleanExternalFill, draw=k8cleanExternalStroke},
  runtimeedge/.style={->, line width=1.05pt, draw=k8cleanRuntime},
  controledge/.style={->, line width=0.95pt, draw=k8cleanControl, dashed},
  artifactedge/.style={->, line width=1.0pt, draw=k8cleanArtifact},
  edgeLabel/.style={font=\scriptsize\sffamily, fill=white, inner sep=2pt, align=center}
]

% Overall cluster and namespace boundaries.
\draw[draw=k8cleanCluster, line width=1.15pt, rounded corners=7pt]
  (0,0.25) rectangle (20,11.6);
\node[font=\bfseries\large\sffamily, anchor=west] at (0.5,11.15)
  {GKE Autopilot Cluster};

\draw[draw=k8cleanAppStroke, line width=1.0pt, rounded corners=6pt]
  (0.5,3.6) rectangle (7.6,10.55);
\node[font=\bfseries\small\sffamily, anchor=west, text=k8cleanAppStroke]
  at (0.8,10.15) {namespace \texttt{philobiblus}};

\draw[draw=k8cleanOpsStroke, line width=1.0pt, rounded corners=6pt]
  (12.25,3.6) rectangle (19.5,10.55);
\node[font=\bfseries\small\sffamily, anchor=west, text=k8cleanOpsStroke]
  at (12.55,10.15) {namespace \texttt{philobiblus-mlops}};

% Application components.
\node[app, minimum width=2.65cm] (gateway) at (2.0,8.55)
  {GKE Gateway /\\HTTPRoute\\\scriptsize API entrypoint};
\node[app, minimum width=2.55cm] (backend) at (5.75,8.55)
  {Backend\\\scriptsize FastAPI, nghiệp vụ};
\node[app, minimum width=2.45cm] (redis) at (2.0,6.05)
  {Redis Cache\\\scriptsize cache, rate limit};
\node[app, minimum width=3.0cm] (recommendation) at (5.75,6.05)
  {Recommendation\\Service\\\scriptsize TF--IDF, :8080};

% MLOps components.
\node[ops, minimum width=2.65cm] (mlflow) at (14.25,8.55)
  {MLflow\\\scriptsize tracking, registry};
\node[ops, minimum width=2.85cm] (trainer) at (17.0,6.05)
  {Retrain CronJob\\\scriptsize train, evaluate, promote};

% Shared resources outside the namespaces.
\node[external, minimum width=3.25cm] (sql) at (3.2,1.45)
  {Cloud SQL PostgreSQL\\\scriptsize application data};
\node[external, minimum width=3.75cm] (gcs) at (10.0,1.45)
  {GCS Artifact Storage\\\scriptsize model, snapshot, release};
\node[external, minimum width=3.25cm] (monitoring) at (16.75,1.45)
  {Cloud Monitoring\\\scriptsize logs, metrics, alerts};

% Runtime flow.
\draw[runtimeedge] (gateway.east) -- (backend.west);
\draw[runtimeedge] (backend.south west) -- (redis.north east);
\node[edgeLabel] at (3.75,7.45) {cache};
\draw[runtimeedge] (backend.south) -- (recommendation.north);
\node[edgeLabel] at (6.2,7.3) {gợi ý};

% MLOps flow.
\draw[controledge] (trainer.north west) -- (mlflow.south east);
\node[edgeLabel] at (15.45,7.55) {run / metric};
\draw[controledge] (trainer.west) -- (recommendation.east);
\node[edgeLabel] at (10.1,6.48) {validation / promotion};

% Data and artifact flow.
\draw[controledge] (backend.south west) -- (3.9,7.75) -- (3.9,2.35) -- (sql.north);
\node[edgeLabel] at (3.55,3.85) {dữ liệu ứng dụng};
\draw[artifactedge] (trainer.south west) -- (gcs.north east);
\node[edgeLabel] at (13.0,3.65) {model / release};
\draw[artifactedge] (gcs.north west) -- (recommendation.south east);
\node[edgeLabel] at (8.5,3.35) {tải model};
\draw[controledge] (trainer.south east) -- (monitoring.north east);
\node[edgeLabel] at (18.15,3.45) {logs / metrics};

% Legend.
\node[anchor=west, font=\scriptsize\sffamily] at (0.7,0.58)
  {\textcolor{k8cleanRuntime}{\rule[0.5ex]{0.65cm}{1pt}} runtime};
\node[anchor=west, font=\scriptsize\sffamily] at (5.3,0.58)
  {\textcolor{k8cleanControl}{\rule[0.5ex]{0.65cm}{1pt}} MLOps control};
\node[anchor=west, font=\scriptsize\sffamily] at (10.8,0.58)
  {\textcolor{k8cleanArtifact}{\rule[0.5ex]{0.65cm}{1pt}} artifact/model};
\end{tikzpicture}%
}

\caption{Các workload và object chính trong cụm GKE của Philobiblus}
\label{fig:kubernetes-architecture}
\end{figure}

\reportsubsection{Triển khai trên Google Cloud Kubernetes Engine}
Helm chart sau khi được kiểm thử và xác minh, được deploy lên GKE Autopilot bằng Terraform. Luồng hoạt động trên Google Cloud gồm:

\textbf{Application traffic:} Người dùng truy cập frontend được phát hành trên GitHub Pages. Request đi qua Cloud Run Proxy rồi tới GKE Gateway và lớp bảo vệ Cloud Armor trước khi vào GKE Autopilot. Trong cluster, namespace ứng dụng tiếp nhận và xử lý request; namespace MLOps không nằm trên đường phục vụ người dùng.

\textbf{Business data:} Backend trong GKE đọc và ghi dữ liệu nghiệp vụ vào Cloud SQL PostgreSQL. Redis vẫn chạy trong namespace ứng dụng để lưu tạm dữ liệu thường dùng và giảm truy vấn lặp lại tới database. Cloud SQL được đặt ngoài cluster nên dữ liệu vẫn nằm trong dịch vụ cơ sở dữ liệu managed, còn workload Kubernetes chỉ sử dụng kết nối được cấp quyền.

\textbf{ML artifacts:} Retrain CronJob và MLflow trong namespace MLOps lưu snapshot, model và thông tin release vào Cloud Storage. Khi cần triển khai phiên bản mới, workload recommendation lấy artifact từ Cloud Storage để nạp model. Artifact Registry lưu các Docker image được build từ source; GKE sử dụng các image này khi tạo hoặc cập nhật workload.

\textbf{Platform \& observability:} Secret Manager cung cấp secret và thông tin nhận dạng cho workload theo cơ chế được cấp quyền, thay vì đưa secret trực tiếp vào chart. Cloud Monitoring và Cloud Logging tiếp nhận log, metric và trạng thái vận hành từ cluster để theo dõi rollout, lỗi và mức sử dụng tài nguyên. Terraform giữ vai trò mô tả và tạo các tài nguyên này theo cùng một cấu hình có thể kiểm tra lại.

Tại thời điểm ngày 04/10/2026, cluster \texttt{philobiblus-dev-gke} đang chạy ổn định, tất cả các tài nguyên được deploy và hoạt động bình thường. 

\begin{figure}[H]
\centering
\definecolor{gcpCleanInk}{HTML}{243447}
\definecolor{gcpCleanProject}{HTML}{64748B}
\definecolor{gcpCleanCluster}{HTML}{2563EB}
\definecolor{gcpCleanClusterFill}{HTML}{EFF6FF}
\definecolor{gcpCleanEdgeStroke}{HTML}{1D4ED8}
\definecolor{gcpCleanEdgeFill}{HTML}{DBEAFE}
\definecolor{gcpCleanOpsStroke}{HTML}{7C3AED}
\definecolor{gcpCleanOpsFill}{HTML}{F5F3FF}
\definecolor{gcpCleanManagedFill}{HTML}{F8FAFC}
\definecolor{gcpCleanManagedStroke}{HTML}{475569}
\definecolor{gcpCleanRuntime}{HTML}{1D4ED8}
\definecolor{gcpCleanData}{HTML}{C2410C}
\definecolor{gcpCleanArtifact}{HTML}{15803D}
\definecolor{gcpCleanInfra}{HTML}{64748B}

\resizebox{\textwidth}{!}{%
\begin{tikzpicture}[
  x=1cm,y=1cm,
  >=latex,
  font=\scriptsize\sffamily,
  every node/.style={text=gcpCleanInk},
  box/.style={draw=gcpCleanInk, line width=0.75pt, rounded corners=3pt, align=center, inner sep=5pt, minimum height=9mm},
  client/.style={box, fill=white, draw=gcpCleanManagedStroke},
  edgebox/.style={box, fill=gcpCleanEdgeFill, draw=gcpCleanEdgeStroke},
  namespace/.style={box, fill=white, draw=gcpCleanCluster, minimum height=1.4cm},
  mlops/.style={box, fill=gcpCleanOpsFill, draw=gcpCleanOpsStroke, minimum height=1.4cm},
  managed/.style={box, fill=gcpCleanManagedFill, draw=gcpCleanManagedStroke, minimum width=3.0cm},
  runtimeedge/.style={->, line width=1.0pt, draw=gcpCleanRuntime},
  dataedge/.style={->, line width=0.9pt, draw=gcpCleanData, dashed},
  artifactedge/.style={->, line width=0.95pt, draw=gcpCleanArtifact},
  infraedge/.style={->, line width=0.85pt, draw=gcpCleanInfra, dashed}
]

% External frontend path.
\node[client, minimum width=1.9cm] (browser) at (1.0,8.7) {Trình duyệt};
\node[client, minimum width=2.35cm] (pages) at (3.75,8.7) {GitHub Pages\\\scriptsize React static site};

% Google Cloud boundary.
\draw[draw=gcpCleanProject, line width=1.15pt, rounded corners=7pt]
  (4.9,0.4) rectangle (20.0,10.4);
\node[font=\bfseries\large\sffamily, anchor=west] at (5.25,9.95)
  {Google Cloud deployment architecture};

\node[edgebox, minimum width=2.55cm] (proxy) at (6.2,8.7)
  {Cloud Run Proxy\\\scriptsize HTTPS entrypoint};
\node[edgebox, minimum width=2.9cm] (gateway) at (9.35,8.7)
  {GKE Gateway + Cloud Armor\\\scriptsize protected API edge};

% Unified GKE cluster.
\draw[draw=gcpCleanCluster, line width=1.15pt, rounded corners=6pt, fill=gcpCleanClusterFill]
  (5.35,1.35) rectangle (12.25,7.65);
\node[font=\bfseries\small\sffamily, anchor=west, text=gcpCleanCluster]
  at (5.7,7.28) {GKE Autopilot Cluster};
\node[namespace, minimum width=2.75cm] (appns) at (7.2,4.5)
  {namespace\\\texttt{philobiblus}\\[2pt]
   \scriptsize Backend\\Redis\\Recommendation};
\node[mlops, minimum width=2.8cm] (mlopsns) at (10.45,4.5)
  {namespace\\\texttt{philobiblus-mlops}\\[2pt]
   \scriptsize Retrain CronJob\\MLflow};

% Managed services outside GKE.
\node[managed] (sql) at (16.4,7.35)
  {Cloud SQL PostgreSQL\\\scriptsize application data};
\node[managed] (gcs) at (16.4,5.95)
  {Cloud Storage\\\scriptsize model, snapshot, release};
\node[managed] (registry) at (16.4,4.55)
  {Artifact Registry\\\scriptsize Docker image};
\node[managed] (secrets) at (16.4,3.15)
  {Secret Manager\\\scriptsize secrets, identity};
\node[managed] (monitoring) at (16.4,1.75)
  {Cloud Monitoring\\\scriptsize logs, metrics, alerts};

% Application traffic.
\draw[runtimeedge] (browser.east) -- (pages.west);
\draw[runtimeedge] (pages.east) -- (proxy.west);
\draw[runtimeedge] (proxy.east) -- (gateway.west);
\draw[runtimeedge] (gateway.south) -- (9.35,7.65);

% GKE-to-managed-service relationships, one horizontal lane per service.
\draw[dataedge] (12.25,7.35) -- (14.9,7.35);
\draw[artifactedge] (12.25,5.95) -- (14.9,5.95);
\draw[infraedge] (12.25,4.55) -- (14.9,4.55);
\draw[infraedge] (12.25,3.15) -- (14.9,3.15);
\draw[infraedge] (12.25,1.75) -- (14.9,1.75);

\node[font=\scriptsize\sffamily, fill=white, inner sep=2pt] at (13.55,7.62) {dữ liệu};
\node[font=\scriptsize\sffamily, fill=white, inner sep=2pt] at (13.55,6.22) {artifact / model};
\node[font=\scriptsize\sffamily, fill=white, inner sep=2pt] at (13.55,4.82) {image};
\node[font=\scriptsize\sffamily, fill=white, inner sep=2pt] at (13.55,3.42) {secret};
\node[font=\scriptsize\sffamily, fill=white, inner sep=2pt] at (13.55,2.02) {logs / metrics};

% Legend.
\node[anchor=west, font=\scriptsize\sffamily] at (5.7,0.9)
  {\textcolor{gcpCleanRuntime}{\rule[0.65ex]{0.65cm}{1pt}} Application traffic};
\node[anchor=west, font=\scriptsize\sffamily] at (9.2,0.9)
  {\textcolor{gcpCleanData}{\rule[0.65ex]{0.65cm}{1pt}} Business data};
\node[anchor=west, font=\scriptsize\sffamily] at (12.4,0.9)
  {\textcolor{gcpCleanArtifact}{\rule[0.65ex]{0.65cm}{1pt}} ML artifacts};
\node[anchor=west, font=\scriptsize\sffamily] at (16.3,0.9)
  {\textcolor{gcpCleanInfra}{\rule[0.65ex]{0.65cm}{1pt}} Platform \& observability};
\end{tikzpicture}%
}

\caption{Luồng triển khai và các resource Google Cloud của Philobiblus}
\label{fig:gcp-architecture}
\end{figure}

\reportsection{c. Ưu điểm và nhược điểm đang gặp phải}
\textbf{Ưu điểm.} TF-IDF có chi phí tính toán thấp, dễ giải thích và vẫn có thể tạo gợi ý khi catalog chưa có nhiều lịch sử tương tác. Mô hình này phù hợp với giai đoạn đầu của Philobiblus vì dữ liệu sách và metadata đã có sẵn, trong khi dữ liệu hành vi người dùng còn hạn chế. Khi model được lưu trong MLflow và artifact được quản lý trên Cloud Storage, quá trình huấn luyện và phát hành có thể truy vết theo snapshot, run và phiên bản model.

Việc triển khai bằng Terraform, Helm và GKE giúp cấu hình hạ tầng và workload có thể kiểm tra, tái tạo và cập nhật theo cùng một quy trình. Backend có thể mở rộng bằng HPA, Redis giảm các truy vấn lặp lại, còn probe và rollout giúp hạn chế việc đưa Pod chưa sẵn sàng vào phục vụ. Cloud Monitoring và Cloud Logging cũng cung cấp dữ liệu để phát hiện lỗi và theo dõi mức sử dụng tài nguyên sau khi triển khai.

\textbf{Nhược điểm và vấn đề đang gặp phải.} Chất lượng gợi ý vẫn phụ thuộc mạnh vào metadata; TF-IDF chưa khai thác được mức độ tương đồng hình thành từ hành vi của cộng đồng người đọc. Các metric của run hiện tại còn có thể ở trạng thái \textit{insufficient data}, vì vậy chưa nên diễn giải điểm đánh giá như một bằng chứng đầy đủ về chất lượng sản phẩm. Khi dữ liệu tương tác đủ lớn, hệ thống cần bổ sung mô hình kết hợp giữa đặc trưng nội dung và hành vi người dùng.

Redis chỉ giữ catalog trong thời gian ngắn và được làm mới khi dữ liệu thay đổi. Nếu Redis tạm thời không khả dụng, Backend bỏ qua cache để tiếp tục phục vụ nhưng sẽ tạo thêm tải cho Cloud SQL. Mỗi Backend Pod dùng pool ba kết nối cơ bản và tối đa thêm hai kết nối, với timeout mười giây; giới hạn này phải được tính cùng HPA để không vượt khả năng của Cloud SQL db-f1-micro.

GKE Autopilot, Cloud SQL, Cloud Storage, Secret Manager và các lớp bảo vệ mạng giúp hệ thống có cấu trúc rõ ràng hơn, nhưng đồng thời làm tăng số lượng tài nguyên phải quản lý, số quyền IAM cần kiểm soát và chi phí vận hành. Cloud Run Proxy bổ sung một lớp trung gian cho frontend, đổi lại tạo thêm một chặng trong luồng truy cập. Đây là các đánh đổi chấp nhận được cho môi trường thực tập và demo, nhưng cần được xem xét lại nếu hệ thống chuyển sang quy mô người dùng lớn hơn.
